\section{Validation}
\label{sec:validation}

The framework is tested against two bodies for which independent gravitational
products exist. Ryugu provides the stronger test, because the Hayabusa2 team
has released per-facet gravitational quantities computed on a shape model that
we can use unchanged, at a stated density and rotation period; the comparison
is therefore facet by facet with every input held fixed. Bennu provides a
complementary test at the level of published summary statistics, and adds a
check that the framework responds correctly to the assumed bulk density.

\subsection{Ryugu: comparison facet by facet}
\label{sec:ryugu}

\citet{kanamaru2019} and the archive accompanying the Hayabusa2 shape analysis
distribute, alongside the $49\,152$-facet SFM shape model of Ryugu, a table of
gravitational potential and gravitational slope for every plate of that model,
computed with a Werner--Scheeres polyhedral solver for a uniform density of
$1200\ \mathrm{kg\,m^{-3}}$ and a rotation period of $7.63$~h. Because shape,
method, density, and spin are all fixed by that specification, any disagreement
must originate in the implementation rather than in the inputs. This is the
strongest test available to us, and the only one in this paper that reaches
below the level of aggregate statistics.

It is worth being precise about what such a test does and does not establish.
The reference products are themselves computed for a homogeneous interior with a
Werner--Scheeres solver, so the comparison is an \emph{implementation}
validation: it confirms that our sign conventions, facet correspondence and
numerical evaluation agree with an independent implementation of the same
physical model. It is not a \emph{physical} validation, and it says nothing
about how well a homogeneous model represents the true gravity field of Ryugu.
That limitation is taken up in Section~\ref{sec:uniformdensity}.

The comparison is run on the full mesh, with no decimation, so that our facets
correspond one-to-one with the archived plates. The correspondence is confirmed
geometrically rather than assumed: matching by index order, the median distance
between our facet centroids and the archived plate centres is $0.04$~m, with a
maximum of $0.1$~m, against a body some $900$~m across. The two meshes are
therefore the same mesh in the same order, and a facet-by-facet difference is
meaningful.

Table~\ref{tab:ryugu} summarises the outcome, and
Figure~\ref{fig:ryuguval} shows the comparison facet by facet. The area-weighted mean slope
agrees to better than $0.01^{\circ}$, and the distribution of per-facet
differences is tight: the RMS difference is $0.48^{\circ}$ and the median
absolute difference $0.34^{\circ}$, with $49\,133$ of the $49\,152$ facets
($99.96\%$) agreeing to within $5^{\circ}$ and $49\,127$ ($99.95\%$) to within
$2^{\circ}$. The Pearson correlation across facets is $0.9989$ for slope and
$0.9898$ for geopotential.

The residual is not uniformly small, and it is worth being explicit about its
tail. Nineteen facets differ by more than $5^{\circ}$ and seven by more than
$10^{\circ}$, the largest by $24.45^{\circ}$. These are a vanishing fraction of
the surface---the $99.9$th percentile of the absolute difference is
$0.95^{\circ}$---but they are real, and Figure~\ref{fig:ryuguval}(b) is plotted
on a logarithmic count axis over the full residual range so that they remain
visible.

\begin{table}[t]
\centering
\small
\caption{Ryugu, facet-by-facet comparison against the archived Hayabusa2
gravitational products. Both fields are computed on the same $49\,152$-facet
shape model at $\dens=1200\ \mathrm{kg\,m^{-3}}$ and $P=7.63$~h.}
\label{tab:ryugu}
\begin{tabular}{lr}
\toprule
Quantity & Value \\
\midrule
Mean slope, this work                       & $15.06^{\circ}$ \\
Mean slope, archived products               & $15.06^{\circ}$ \\
Mean difference                             & $-0.00^{\circ}$ \\
RMS difference                              & $0.48^{\circ}$ \\
Median absolute difference                  & $0.34^{\circ}$ \\
$99.9$th percentile of $\lvert$difference$\rvert$ & $0.95^{\circ}$ \\
Maximum absolute difference                 & $24.45^{\circ}$ \\
Facets within $2^{\circ}$                   & $49\,127$ \ ($99.95\%$) \\
Facets within $5^{\circ}$                   & $49\,133$ \ ($99.96\%$) \\
Facets differing by more than $5^{\circ}$   & $19$ \\
Facets differing by more than $10^{\circ}$  & $7$ \\
Pearson $r$, slope                          & $0.9989$ \\
Pearson $r$, geopotential                   & $0.9898$ \\
Pearson $r$, slope on facets below $15^{\circ}$ & $0.9936$ \\
\bottomrule
\end{tabular}
\end{table}

Two features of the residual are worth noting. First, removing the mean bias
does not reduce the RMS difference, which stays at $0.48^{\circ}$: the
disagreement is scatter rather than offset, as expected if it arises from
differences in how the two implementations evaluate the field near the surface
rather than from a systematic error in either. Second, the correlation is just
as high on the flattest part of the surface---$0.9936$ restricted to facets
whose archived slope is below $15^{\circ}$---as it is overall. A sign or
convention error would degrade the flat facets preferentially, because there
the two competing contributions are most nearly balanced; that the flat facets
agree as well as the steep ones is a useful independent indication that the
conventions of Section~\ref{sec:signs} are correct.

The agreement also holds as a function of position on the body. Binning both
data sets by latitude gives the profile in Table~\ref{tab:ryugulat}: a maximum
near the equator, a minimum at intermediate latitude, and a rise toward the
poles. All four bands agree to better than $0.7^{\circ}$, and the two lower
bands to better than $0.2^{\circ}$. The equatorial maximum and mid-latitude
minimum are the signature of the equatorial ridge of a top-shaped body, and the
framework reproduces both without adjustment.

\begin{table}[t]
\centering
\small
\caption{Ryugu, area-weighted mean slope by latitude band, compared with the
archived products.}
\label{tab:ryugulat}
\begin{tabular}{lrrr}
\toprule
$\lvert\text{latitude}\rvert$ & This work & Archive & Difference \\
\midrule
$0$--$20^{\circ}$   & $15.12^{\circ}$ & $15.20^{\circ}$ & $-0.08^{\circ}$ \\
$20$--$40^{\circ}$  & $12.23^{\circ}$ & $12.40^{\circ}$ & $-0.17^{\circ}$ \\
$40$--$60^{\circ}$  & $15.66^{\circ}$ & $15.00^{\circ}$ & $+0.66^{\circ}$ \\
$60$--$90^{\circ}$  & $21.25^{\circ}$ & $20.80^{\circ}$ & $+0.45^{\circ}$ \\
\midrule
Global              & $15.06^{\circ}$ & $15.06^{\circ}$ & $-0.00^{\circ}$ \\
\bottomrule
\end{tabular}
\end{table}

\begin{figure}[tp]
\centering
\includegraphics[width=\textwidth]{fig1.pdf}
\caption{Per-facet validation of the computed slope field against the
gravitational products released with the Hayabusa2 shape analysis of Ryugu,
on the identical $49\,152$-facet model at $\dens = 1200\ \mathrm{kg\,m^{-3}}$
and $P = 7.63$~h. (a) Facet-by-facet comparison; the colour scale is
logarithmic in the number of facets per bin, and the dashed line is the $1{:}1$
relation. (b) Distribution of the per-facet difference. The agreement is
scatter rather than offset: removing the mean bias leaves the RMS difference
unchanged.}
\label{fig:ryuguval}
\end{figure}

\subsection{Bennu: global structure}
\label{sec:bennu}

For Bennu the OLA shape model \citep{barnouin2020,daly2020} is decimated to
$30\,561$ facets and evaluated at the measured bulk density
$\dens=1.194\ \mathrm{g\,cm^{-3}}$ and the rotation period $P=4.2961$~h. The
resulting area-weighted mean slope is $15.45^{\circ}$, with a median of
$14.94^{\circ}$; the distribution is concentrated between $10^{\circ}$ and
$20^{\circ}$ (which together hold $48.7\%$ of the surface area), only $3.9\%$
of facets exceed $30^{\circ}$, and just two facets out of $30\,561$ exceed
$90^{\circ}$. These figures sit within the range of global mean slopes reported
from OSIRIS-REx analyses \citep{scheeres2019,jawin2020}.

The latitudinal structure is more diagnostic than the global mean, because it
is set by the competition between the equatorial ridge and the centrifugal
term. The profile, given in Table~\ref{tab:bennulat}, rises from $10.45^{\circ}$
within $10^{\circ}$ of the equator to a broad maximum of $19.89^{\circ}$ at
$45$--$60^{\circ}$ latitude before falling to $17.83^{\circ}$ near the pole.
This reproduces the pattern reported from mission data---low slopes across the
equatorial region, a transition near $\pm 20^{\circ}$, and higher slopes at
mid to high latitude \citep{scheeres2019}---both in shape and, to within a
degree or two, in value; Figure~\ref{fig:bennuval}(a) shows the profile against
the published values.

\begin{table}[t]
\centering
\small
\caption{Bennu, area-weighted mean slope by latitude band at
$\dens = 1.194\ \mathrm{g\,cm^{-3}}$ and $\Nf = 30\,561$.}
\label{tab:bennulat}
\begin{tabular}{lrrr}
\toprule
$\lvert\text{latitude}\rvert$ & Mean & Median & Facets \\
\midrule
$0$--$10^{\circ}$   & $10.45^{\circ}$ & $9.70^{\circ}$  & $6202$ \\
$10$--$20^{\circ}$  & $12.02^{\circ}$ & $11.36^{\circ}$ & $5804$ \\
$20$--$30^{\circ}$  & $16.06^{\circ}$ & $15.48^{\circ}$ & $4949$ \\
$30$--$45^{\circ}$  & $19.10^{\circ}$ & $18.43^{\circ}$ & $5856$ \\
$45$--$60^{\circ}$  & $19.89^{\circ}$ & $19.32^{\circ}$ & $4230$ \\
$60$--$90^{\circ}$  & $17.83^{\circ}$ & $17.36^{\circ}$ & $3520$ \\
\bottomrule
\end{tabular}
\end{table}

The geopotential of Eq.~\eqref{eq:geopot} carries the same information in a
form directly relevant to regolith transport. Its minimum lies at the equator
and it increases toward the poles, so that loose material on the surface is
driven equatorward; this is the geopotential structure inferred for Bennu from
mission data and invoked to explain its equatorial accumulation
\citep{scheeres2019,jawin2020}. The framework recovers it from the shape model, the
assumed density and the specified rotation period.

\subsection{Bennu: response to the assumed density}
\label{sec:densitysweep}

A shape-based calculation must assume a bulk density, and it is worth
establishing that the framework responds to that assumption in the physically
correct way rather than merely producing a plausible number at one value. For a
rotating body the relevant control parameter is the ratio of centrifugal to
gravitational acceleration; since the centrifugal term is fixed by the spin and
the shape, raising the assumed density strengthens self-gravity and reduces
that ratio, which must lower the slope.

Because the field is linear in $\dens$ (Section~\ref{sec:solver}), the entire
sweep costs one solver evaluation. Table~\ref{tab:sweep} gives the result, and
Figure~\ref{fig:bennuval}(b) displays it. The
mean slope falls from $25.55^{\circ}$ at $\dens = 0.85\ \mathrm{g\,cm^{-3}}$ to
$15.45^{\circ}$ at the measured $1.194\ \mathrm{g\,cm^{-3}}$, while the
equatorial centrifugal-to-gravitational ratio falls from $0.751$ to $0.534$.
The comparison values reported from OSIRIS-REx analyses for the same exercise
are approximately $24^{\circ}$ at $0.85$ and $15^{\circ}$ at
$1.15\ \mathrm{g\,cm^{-3}}$ \citep{scheeres2019}; the framework returns
$25.55^{\circ}$ and $16.06^{\circ}$, agreeing to $1.5^{\circ}$ and
$1.1^{\circ}$ respectively. The agreement is therefore quantitative across the
range, not merely in trend, and at the measured density the equatorial and
polar means---$11.19^{\circ}$ and $17.81^{\circ}$---bracket the published
equatorial and high-latitude figures of roughly $12^{\circ}$ and $18^{\circ}$.

\begin{table}[t]
\centering
\small
\caption{Bennu, area-weighted mean slope as a function of the assumed bulk
density at $\Nf = 30\,561$. Here $\text{eq}$ denotes
$\lvert\text{lat}\rvert<20^{\circ}$, $\text{mid}$ denotes
$20$--$40^{\circ}$, and $\text{pol}$ denotes $>60^{\circ}$;
$\lVert\acf\rVert/\lVert\ngrav\rVert$ is the equatorial mean.}
\label{tab:sweep}
\begin{tabular}{rrrrrr}
\toprule
$\dens$ & Mean & eq & mid & pol & $\lVert\acf\rVert/\lVert\ngrav\rVert$ \\
\midrule
$0.850$ & $25.55^{\circ}$ & $23.07^{\circ}$ & $29.38^{\circ}$ & $21.28^{\circ}$ & $0.751$ \\
$1.000$ & $19.18^{\circ}$ & $14.57^{\circ}$ & $22.49^{\circ}$ & $19.43^{\circ}$ & $0.638$ \\
$1.150$ & $16.06^{\circ}$ & $11.61^{\circ}$ & $18.19^{\circ}$ & $18.12^{\circ}$ & $0.555$ \\
$1.194$ & $15.45^{\circ}$ & $11.19^{\circ}$ & $17.26^{\circ}$ & $17.81^{\circ}$ & $0.534$ \\
\bottomrule
\end{tabular}
\end{table}

\begin{figure}[tp]
\centering
\includegraphics[width=\textwidth]{fig2.pdf}
\caption{Bennu. (a) Area-weighted mean and median slope by latitude band at the
measured bulk density; shaded bands mark the equatorial and high-latitude values
reported from OSIRIS-REx analyses. The profile reproduces the observed pattern
of low equatorial slopes rising to a mid-latitude maximum. (b) Dependence of the
slope on the assumed bulk density, resolved by latitude zone. Stars mark the
published global means at $0.85$ and $1.15\ \mathrm{g\,cm^{-3}}$. The
equatorial zone responds most strongly, as expected for a centrifugal effect
that scales with distance from the spin axis.}
\label{fig:bennuval}
\end{figure}

The equatorial band responds most strongly, falling by nearly $12^{\circ}$
across the range while the polar band moves by only $3.5^{\circ}$. That is the
expected signature: the centrifugal term acts perpendicular to the spin axis
and so has its largest effect where the surface is farthest from that axis.
The framework therefore reproduces not only the magnitude of the
density dependence but its spatial distribution.

\subsection{Summary}
\label{sec:valsummary}

Taken together, the two tests bound the accuracy of the framework from
different directions. Ryugu establishes that, given identical inputs, the
implementation reproduces an independent Werner--Scheeres calculation facet by
facet, with an RMS difference of half a degree and $99.96\%$ of facets
agreeing to within $5^{\circ}$. Bennu establishes that, given a shape model, an
assumed density and a specified spin state, the framework recovers the
published global mean, the
latitudinal structure, the equatorial geopotential minimum, and the
quantitative dependence on assumed density.

Neither test says anything about how these numbers would change had a different
mesh been used, and that question is the subject of the next two sections.
